\documentclass[11pt,a4paper]{article}

\usepackage[utf8]{inputenc}
\usepackage[T1]{fontenc}
\usepackage{lmodern}
\usepackage[brazilian]{babel}
\usepackage[a4paper,margin=2.5cm]{geometry}
\usepackage{microtype}
\usepackage[table]{xcolor}
\usepackage{booktabs}
\usepackage{array}
\usepackage{enumitem}
\usepackage{titlesec}
\usepackage{caption}
\usepackage[section]{placeins}
\usepackage{tikz}
\usetikzlibrary{arrows.meta,positioning,calc,decorations.pathreplacing}
\usepackage[colorlinks=true,linkcolor=azul,urlcolor=azul]{hyperref}

\definecolor{azul}{RGB}{31,78,121}
\definecolor{azulclaro}{RGB}{222,234,247}
\definecolor{laranja}{RGB}{191,87,0}
\definecolor{laranjaclaro}{RGB}{255,236,214}
\definecolor{verde}{RGB}{46,125,50}
\definecolor{verdeclaro}{RGB}{226,242,227}
\definecolor{cinza}{RGB}{150,150,150}
\definecolor{cinzaclaro}{RGB}{238,238,238}
\definecolor{cinzaescuro}{RGB}{80,80,80}

\titleformat{\section}{\large\bfseries\color{azul}}{\thesection}{0.8em}{}
\titleformat{\subsection}{\normalsize\bfseries}{\thesubsection}{0.8em}{}
\captionsetup{font=small,labelfont=bf}

\setlist{itemsep=2pt,topsep=4pt}
\setlength{\emergencystretch}{3em}
\newcommand{\cod}[1]{\texttt{#1}}

\tikzset{
  seta/.style={-{Stealth[length=2.2mm]}, thick, draw=azul},
  caixa/.style={draw=azul, rounded corners=2pt, fill=azulclaro, text width=7.6cm, align=center, inner sep=5pt,
    font=\small},
  futura/.style={caixa, dashed, draw=cinza, fill=white, text=cinzaescuro},
  rotulo/.style={font=\scriptsize\ttfamily, text=cinzaescuro, align=right},
}

\hypersetup{pdftitle={Relatório técnico: apoio computacional ao projeto de tese},
  pdfauthor={Caio Victor Ferreira do Nascimento}}

\begin{document}

\begin{titlepage}
\centering
{\large UNIVERSIDADE FEDERAL DO PIAUÍ\\Departamento de Computação\par}
\vspace{4cm}
{\LARGE\bfseries\color{azul} Relatório Técnico\par}
\vspace{0.8cm}
{\large Trabalho de apoio computacional ao projeto de tese\\[0.3cm]
\emph{Desenvolvimento de modelos de inteligência artificial para predição de sepse tardia em Unidade de Terapia
Intensiva Neonatal}\par}
\vspace{2.5cm}
\begin{tabular}{rl}
Aluno: & Caio Victor Ferreira do Nascimento\\
Professor: & Prof. Dr. Vinícius Ponte Machado\\
\end{tabular}
\vfill
{\large Teresina, 31 de outubro de 2026\par}
\end{titlepage}

% ==================================================================
\section{Objetivo}

Este relatório resume o que foi desenvolvido até agora no pipeline de aprendizado de máquina do projeto de tese. Ele
lista os requisitos que o projeto define, mostra como cada um foi implementado e indica o que ainda falta. As páginas
citadas são as do documento do projeto.

O trabalho concentrou-se na parte do pipeline que \textbf{não depende dos dados}: algoritmos, hiperparâmetros,
tratamento do desbalanceamento, validação cruzada, avaliação e explicabilidade. O código está no repositório, no ramo
\cod{create-pipeline}, e tem 79 testes automáticos, todos aprovados.

% ==================================================================
\section{Requisitos implementados}

A Figura~\ref{fig:fluxo} mostra o caminho dos dados no pipeline. Cada caixa traz o arquivo responsável, e cada seta, o
que esse arquivo grava. A caixa tracejada, dos dados reais, ainda não existe.

\begin{figure}[ht]
\centering
\begin{tikzpicture}[node distance=9mm]
\node[futura] (dados) {\textbf{Dados reais}\\extração do prontuário SOUL MV (a implementar)};
\node[caixa, below=of dados] (div) {\textbf{divisao.py}\\divisão 80/20, estratificada pelo desfecho};
\node[caixa, below=14mm of div] (val) {\textbf{validacao.py} e \textbf{modelos.py}\\
  pré-processamento (a definir)\\
  8 algoritmos e 3 modos de desbalanceamento\\
  GridSearchCV e validação cruzada aninhada (5 dobras)};
\node[caixa, below=14mm of val] (ava) {\textbf{avaliacao.py}\\escolhe o vencedor e só então abre o teste};
\node[caixa, below=14mm of ava] (exp) {\textbf{explicacao.py}\\valores SHAP do vencedor e modelo final};
\node[rotulo, align=center, below=3mm of exp] {grava modelo\_final.joblib e .json\\e a importância SHAP};
\draw[seta, dashed] (dados) -- (div);
\draw[seta] (div) -- node[left=2mm, rotulo] {treino.csv} (val);
\draw[seta] (val) -- node[left=2mm, rotulo] {comparacao\_*.csv, oof\_*.csv\\modelos\_*/ (um .joblib por algoritmo)} (ava);
\draw[seta] (ava) -- node[left=2mm, rotulo] {vencedor\_*.json, curvas ROC\\matriz de confusão} (exp);
\draw[seta, draw=laranja, dashed] (div.east) -- ++(1.0,0) |- node[pos=0.25, right, rotulo, align=left, text=laranja]
  {teste.csv\\lacrado até\\a avaliação} (ava.east);
\end{tikzpicture}
\caption{Fluxograma do pipeline, com o arquivo responsável por cada etapa e o que cada um grava.}
\label{fig:fluxo}
\end{figure}

\subsection{Oito algoritmos de classificação}

O Quadro~4 (p.~44) lista oito algoritmos: Regressão Logística (RL), Naive Bayes (NB), K-vizinhos mais próximos (KNN),
Árvore de Decisão (DT), Random Forest (RF), Gradient Boosting (GB), XGBoost e rede neural MLP. Todos estão no arquivo
\cod{modelos.py}, que funciona como um catálogo. Para cada algoritmo, ele define o estimador, os parâmetros fixos e as
grades de hiperparâmetros. Sete vêm da biblioteca scikit-learn, e o XGBoost vem da biblioteca xgboost.

Cada algoritmo é montado num \emph{pipeline} com um espaço reservado para o pré-processamento, sempre antes do modelo.
Esse espaço fica vazio até os dados reais chegarem.

\subsection{Busca de hiperparâmetros}

O projeto prevê a busca com \cod{GridSearchCV} (p.~45). Cada algoritmo tem uma grade enxuta, para os testes, e uma
grade completa, prevista para os dados reais, com 59 e 325 combinações no total. Testes automáticos treinam cada
algoritmo com a primeira, a do meio e a última combinação de cada grade. Assim, um hiperparâmetro inválido aparece
antes de uma busca longa, e não no meio dela.

\subsection{Divisão entre treino e teste}

O arquivo \cod{divisao.py} lê a tabela, confere a coluna de identificação e o desfecho 0/1 e separa 80\% dos pacientes
para treino e 20\% para teste (p.~45). A divisão é estratificada: as duas partes mantêm a mesma proporção de sepse. O
teste fica lacrado até a avaliação final.

\subsection{Validação cruzada com cinco dobras}

O arquivo \cod{validacao.py} faz a validação cruzada com cinco dobras (p.~45), usando \cod{StratifiedKFold}. A busca
de hiperparâmetros escolhe o melhor modelo justamente pelo desempenho nas dobras, então a nota dela é otimista. Por
isso, a comparação entre algoritmos usa a validação aninhada da Figura~\ref{fig:dobras}: a busca inteira acontece
dentro de cada rodada, e o modelo escolhido é medido numa dobra que ele nunca viu.

\begin{figure}[ht]
\centering
\begin{tikzpicture}[font=\small]
  % validação externa
  \node[anchor=south west, font=\small\bfseries] at (0,0.05) {Validação externa};
  \foreach \r in {1,...,5} {
    \node[anchor=east] at (-0.1,{-0.7*\r+0.225}) {Rodada \r};
    \foreach \d in {1,...,5} {
      \ifnum\d=\r \def\cor{laranjaclaro}\def\borda{laranja}\else \def\cor{azulclaro}\def\borda{azul}\fi
      \filldraw[fill=\cor, draw=\borda] ({\d-1},{-0.7*\r}) rectangle ++(0.95,0.45);
    }
  }
  % destaque da rodada 1 e ligação com a busca interna
  \draw[dashed, thick, draw=azul] (0.97,-0.78) rectangle (5.0,-0.17);
  \draw[seta] (5.05,-0.47) to[out=0,in=180] (6.35,-1.87);
  % busca interna
  \node[anchor=south west, font=\small\bfseries] at (6.5,0.05) {Busca interna da rodada 1};
  \foreach \r in {1,...,5} {
    \foreach \d in {1,...,5} {
      \ifnum\d=\r \def\cor{verdeclaro}\def\borda{verde}\else \def\cor{azulclaro}\def\borda{azul}\fi
      \filldraw[fill=\cor, draw=\borda] ({6.5+0.8*(\d-1)},{-0.7*\r}) rectangle ++(0.76,0.45);
    }
  }
  % legenda
  \begin{scope}[shift={(0,-4.4)}]
    \filldraw[fill=azulclaro, draw=azul] (0,0) rectangle ++(0.4,0.3);
    \node[anchor=west] at (0.5,0.15) {treino};
    \filldraw[fill=verdeclaro, draw=verde] (2.2,0) rectangle ++(0.4,0.3);
    \node[anchor=west] at (2.7,0.15) {escolha dos hiperparâmetros};
    \filldraw[fill=laranjaclaro, draw=laranja] (7.6,0) rectangle ++(0.4,0.3);
    \node[anchor=west] at (8.1,0.15) {medida do modelo escolhido};
  \end{scope}
\end{tikzpicture}
\caption{Validação cruzada aninhada. Em cada rodada externa, os quatro quintos de treino são redivididos em cinco
dobras para escolher os hiperparâmetros, e o modelo escolhido é medido no quinto restante. O conjunto de teste fica
fora de todo o processo.}
\label{fig:dobras}
\end{figure}

\subsection{Tratamento do desbalanceamento}

O projeto prevê pesos de classe e reamostragem sintética (p.~46), porque os casos de sepse são minoria. São três modos,
comparados lado a lado, e o Quadro~\ref{tab:desb} mostra como cada algoritmo recebe o ajuste. O GB e o XGBoost
precisaram de uma classe própria. O XGBoost, por exemplo, só aceita um peso fixo, que, calculado com o treino inteiro,
incluiria a dobra de validação. A classe própria recalcula o peso a cada treino.

\begin{table}[ht]
\centering
\small
\caption{Como cada algoritmo recebe o ajuste em cada modo de desbalanceamento.}
\label{tab:desb}
\begin{tabular}{@{}lccc@{}}
\toprule
Algoritmo & Sem ajuste & Pesos de classe & SMOTE \\
\midrule
RL      & --- & \cellcolor{verdeclaro}nativo & \cellcolor{laranjaclaro}SMOTE \\
NB      & --- & \cellcolor{cinzaclaro}sem suporte & \cellcolor{laranjaclaro}SMOTE \\
KNN     & --- & \cellcolor{cinzaclaro}sem suporte & \cellcolor{laranjaclaro}SMOTE \\
DT      & --- & \cellcolor{verdeclaro}nativo & \cellcolor{laranjaclaro}SMOTE \\
RF      & --- & \cellcolor{verdeclaro}nativo & \cellcolor{laranjaclaro}SMOTE \\
GB      & --- & \cellcolor{azulclaro}classe própria & \cellcolor{laranjaclaro}SMOTE \\
XGBoost & --- & \cellcolor{azulclaro}classe própria & \cellcolor{laranjaclaro}SMOTE \\
MLP     & --- & \cellcolor{cinzaclaro}sem suporte & \cellcolor{laranjaclaro}SMOTE \\
\bottomrule
\end{tabular}
\end{table}

O SMOTE cria casos de sepse por interpolação entre casos vizinhos e, como age nos dados, cobre os oito algoritmos. Ele
roda só nas dobras de treino, como mostra a Figura~\ref{fig:smote}. Um teste automático confirma que, se ele rodasse
antes da divisão, o desempenho ficaria artificialmente alto.

\begin{figure}[ht]
\centering
\begin{tikzpicture}[node distance=8mm, every node/.style={font=\small}]
\node[caixa, text width=3.4cm] (tr) {dobras de treino};
\node[caixa, text width=3.6cm, fill=laranjaclaro, draw=laranja, right=of tr] (sm) {SMOTE\\iguala as classes};
\node[caixa, text width=3.0cm, right=of sm] (aj) {treino do modelo};
\node[caixa, text width=3.4cm, below=10mm of tr] (va) {dobra de validação};
\node[caixa, text width=3.0cm, below=10mm of aj] (pr) {previsão e métricas};
\draw[seta] (tr) -- (sm);
\draw[seta] (sm) -- (aj);
\draw[seta] (aj) -- (pr);
\draw[seta] (va) -- node[above, font=\scriptsize, text=cinzaescuro] {sem SMOTE: casos reais apenas} (pr);
\end{tikzpicture}
\caption{O SMOTE age só no treino de cada dobra. A dobra de validação é avaliada com os casos originais.}
\label{fig:smote}
\end{figure}

\subsection{Métricas de desempenho}

As métricas do Quadro~5 (p.~45) estão em \cod{validacao.py}: F1-score, acurácia, precisão, sensibilidade e AUROC,
além da especificidade, citada no Resumo do projeto. O F1, com limiar 0,5, escolhe o melhor modelo, e a AUROC desempata.
Os testes conferem as métricas com contas feitas à mão.

\subsection{Comparação em tabela e curvas ROC}

O arquivo \cod{avaliacao.py} escolhe o vencedor entre as 24 combinações de algoritmo e modo pela validação aninhada e só
depois abre o teste. Em seguida, mede todas as combinações no teste e calcula o intervalo de confiança de 95\% do
vencedor por \emph{bootstrap}. Gera também a tabela comparativa, as curvas ROC e a matriz de confusão pedidas na
p.~46. Ainda não foi executado com dados reais.

\subsection{Explicabilidade com SHAP}

O arquivo \cod{explicacao.py} calcula os valores SHAP do vencedor (p.~46), isto é, a contribuição de cada variável para
a probabilidade prevista de cada paciente. Ele também salva o modelo final com a lista de colunas esperadas e tem a
função \cod{prever}, que devolve a probabilidade de sepse de novos pacientes. Ainda não foi executado com dados
reais.

% ==================================================================
\section{Por que os dados ainda não foram trabalhados}

Os dados reais ainda não existem. A coleta depende da aprovação nos Comitês de Ética e do acesso à API do prontuário
SOUL MV e está prevista para 2027 (p.~48). Além disso, as dúvidas da seção~\ref{sec:duvidas} mudam quais dados devem
ser coletados e como. Por isso, o pré-processamento não foi implementado, e não há dados no repositório.

Os testes automáticos usam uma tabela numérica genérica, criada na memória durante o teste e sem nenhum significado
clínico. Ela serve só para verificar que os algoritmos treinam e que nada vaza entre as dobras. Nenhum número dela é
resultado do estudo.

% ==================================================================
\section{Dúvidas sobre o projeto}\label{sec:duvidas}

A Figura~\ref{fig:linha} situa as duas primeiras dúvidas na linha do tempo de uma internação.

\begin{figure}[ht]
\centering
\begin{tikzpicture}[font=\scriptsize]
  % ----- recém-nascido investigado
  \node[anchor=west, font=\small\bfseries] at (0,3.3) {Recém-nascido com suspeita de sepse};
  \fill[laranjaclaro] (8.6,-0.15) rectangle (14,0.15);
  \draw[thick, draw=cinzaescuro, -{Stealth[length=2mm]}] (0,0) -- (14.2,0);
  \foreach \x/\t in {0/admissão, 2.2/72 horas, 6/suspeita\\(hemocultura), 8.6/diagnóstico\\(CID), 13/última\\evolução} {
    \draw[thick] (\x,-0.15) -- (\x,0.15);
    \node[below, align=center] at (\x,-0.2) {\t};
  }
  % janela dos exames
  \draw[decorate, decoration={brace, amplitude=4pt}, thick, draw=azul] (4.8,0.3) -- (7.3,0.3)
    node[midway, above=4pt, align=center, text=azul] {exames: última medida antes\\e até 24 h depois (p. 39--40)};
  % totais da internação
  \draw[thick, draw=laranja, {Stealth[length=2mm]}-{Stealth[length=2mm]}] (0,2.3) -- (13,2.3)
    node[midway, above, text=laranja] {totais da internação: dias de UTIN, ventilação e cateter (Quadro 2)};
  % última evolução
  \draw[seta, draw=laranja] (13,1.9) -- (13,0.25);
  \node[anchor=east, text=laranja, align=right] at (12.9,1.5) {variáveis da\\última evolução (p. 38)};
  \node[text=laranja] at (10.8,0.4) {depois do diagnóstico};
  % ----- recém-nascido sem suspeita
  \node[anchor=west, font=\small\bfseries] at (0,-2.0) {Recém-nascido sem suspeita de sepse};
  \draw[thick, draw=cinzaescuro, -{Stealth[length=2mm]}] (0,-3.3) -- (14.2,-3.3);
  \foreach \x/\t in {0/admissão, 2.2/72 horas, 13/última\\evolução} {
    \draw[thick] (\x,-3.45) -- (\x,-3.15);
    \node[below, align=center] at (\x,-3.5) {\t};
  }
  \node[circle, draw=laranja, thick, text=laranja, font=\small\bfseries, inner sep=1pt] at (6,-3.3) {?};
  \node[below, align=center, text=laranja] at (6,-3.55) {marco dos exames\\não definido};
\end{tikzpicture}
\caption{Onde o projeto manda colher as variáveis. A faixa laranja marca o período depois do diagnóstico. Os totais da
internação e a última evolução incluem esse período, e, para quem não teve suspeita, não há marco para os exames.}
\label{fig:linha}
\end{figure}

\begin{enumerate}
\item \textbf{Qual pergunta o modelo responde.} A introdução fala em prever a sepse com antecedência, entre todos os
internados (p.~13). Já os exames são colhidos no momento da suspeita (p.~39 e 40), que não existe para quem nunca foi
investigado. Se casos e controles tiverem exames de momentos diferentes, o modelo pode aprender a reconhecer quem
estava sendo investigado, e não a sepse.
\item \textbf{Variáveis que podem ser consequência da sepse.} A última evolução e os totais da internação vêm depois do
diagnóstico (p.~38). O resultado da hemocultura e o tempo até o diagnóstico também podem entregar o desfecho ao modelo.
\item \textbf{Excluir ou imputar dados ausentes.} A p.~38 manda excluir prontuários incompletos, e a p.~40 manda
imputar os valores ausentes.
\item \textbf{Escores clínicos.} qSOFA e NEWS2 usam pontos de corte de adultos, que quase todo recém-nascido atinge
(p.~33). O nSOFA, feito para neonatos, é citado no texto, mas não está na ficha de coleta.
\item \textbf{Critério de sucesso.} O projeto não define o que é ``desempenho adequado'' (p.~14) nem o limiar de
decisão do modelo.
\end{enumerate}

% ==================================================================
\section{O que falta implementar}

\begin{itemize}
\item Extração dos dados do SOUL MV pela API e cálculo dos escores clínicos (p.~43).
\item Pré-processamento, que depende dos dados reais e das dúvidas acima.
\item Execução do pipeline completo com os dados reais, incluindo a avaliação final e o SHAP.
\item SVM e redes profundas, citados no Quadro~4, com a regra de exclusão abaixo de 70\%.
\item Análise das variáveis associadas à sepse (objetivo específico~1).
\end{itemize}

% ==================================================================
\section{Alterações quando os dados chegarem}

\begin{itemize}
\item Conferir os nomes das colunas de identificação e de desfecho, hoje os do dicionário de variáveis do projeto.
\item Implementar o pré-processamento e colocá-lo no espaço reservado do pipeline.
\item Excluir ou recortar no tempo as variáveis que dependerem das respostas às dúvidas 1 e 2.
\item Rodar a grade completa de hiperparâmetros e revisar seus valores.
\end{itemize}

Os algoritmos, a validação cruzada, o tratamento do desbalanceamento, a avaliação e o SHAP não precisam mudar.

\end{document}
